% Источники: exam/materials/преподаватель/2020/06-04-2020-Model L7(2).pdf; exam/materials/преподаватель/2020/06-04-2020-Model L7(3).pdf; exam/materials/моделирование.pdf, раздел 39; GitHub HumsterProgrammer/iu9-notes/8sem/Моделирование/Моделирование_лекция-13_26-03-26.md; GitHub HumsterProgrammer/iu9-notes/8sem/Моделирование/Моделирование_лаб-03_23-03-26.md.
\section{Кластерный анализ в математическом моделировании.}

Решение задачи кластеризации принципиально неоднозначно, и тому есть несколько причин:
\begin{itemize}
  \item не существует однозначно наилучшего критерия качества кластеризации;
  \item известен целый ряд эвристических критериев, а также ряд алгоритмов, не имеющих чётко выраженного критерия, но осуществляющих достаточно разумную кластеризацию «по построению»;
  \item число кластеров, как правило, неизвестно заранее и устанавливается в соответствии с некоторым субъективным критерием;
  \item различия в стандартизации переменных;
  \item результат кластеризации существенно зависит от метрики, выбор которой, как правило, также субъективен и определяется экспертом.
\end{itemize}

Различают агломеративные (иерархические) и дивизимные (неиерархические) методы.

\textbf{Агломеративные методы} характеризуются последовательным объединением исходных элементов и соответствующим уменьшением числа кластеров. В начале работы алгоритма все объекты являются отдельными кластерами. На первом шаге наиболее похожие объекты объединяются в кластер. На последующих шагах объединение продолжается до тех пор, пока все объекты не будут составлять один кластер.

При большом количестве наблюдений иерархические методы кластерного анализа не пригодны. В таких случаях используют неиерархические методы, основанные на разделении, которые представляют собой итеративные методы дробления исходной совокупности.

\textbf{Дивизимные методы} являются логической противоположностью агломеративным методам. В начале работы алгоритма все объекты принадлежат одному кластеру, который на последующих шагах делится на меньшие кластеры.

Иерархические алгоритмы связаны с построением дендрограмм, которые являются результатом иерархического кластерного анализа. Дендрограмма описывает близость отдельных точек и кластеров друг к другу, представляет в графическом виде последовательность объединения (разделения) кластеров.

\textbf{Методы:}
\begin{itemize}
  \item \textbf{$k$-средних.} Алгоритм $k$-средних предполагает построение $k$ кластеров, расположенных на возможно больших расстояниях друг от друга. Общая идея алгоритма: заданное фиксированное число $k$ кластеров, наблюдения сопоставляются кластерам так, что средние в кластере (для всех переменных) максимально возможно отличаются друг от друга.
  \item \textbf{Метод Варда.} Метод Варда является иерархическим агломеративным методом. В качестве расстояния между кластерами берется прирост суммы квадратов расстояний объектов до центров кластеров, получаемый в результате их объединения.
  \item \textbf{Метод ближайшего соседа.} Расстояние между двумя кластерами определяется расстоянием между двумя наиболее близкими объектами (ближайшими соседями) в различных кластерах.
  \item \textbf{Метод наиболее удаленного соседа.} То же самое, но самые удаленные.
\end{itemize}

Метрики любые, наиболее популярные $L_p$ ($p$ — натуральное число):
\[
  L_p = \left(\sum_i |x_i-y_i|^p\right)^{1/p}.
\]

\clearpage
